\documentclass[12pt,a4paper]{article}
\usepackage[utf8]{inputenc}
\usepackage[indonesian]{babel}
\usepackage{amsmath}
\usepackage{amsfonts}
\usepackage{amssymb}
\usepackage{geometry}
\usepackage{hyperref}

\geometry{margin=2.5cm}

\title{Metode Secant dalam Komputasi Numerik}
\author{Ahmad Fakhrul Bawani}
\date{}

\begin{document}

\maketitle

\section{Pendahuluan}
Metode Secant merupakan perbaikan dari metode Newton-Raphson bagi fungsi yang turunannya sulit untuk dicari. Dalam metode ini, turunan pertama $f'(x)$ didekati dengan kemiringan garis (tali busur atau \textit{secant}) yang melalui dua titik hampiran awal. Berbeda dengan Biseksi, metode ini termasuk dalam kategori metode terbuka karena tidak selalu membutuhkan akar di dalam rentang awal.

\section{Konsep dan Penurunan Rumus}
Ide utama metode Secant adalah mengganti turunan $f'(x_n)$ pada rumus Newton-Raphson:
\begin{equation}
f'(x_n) \approx \frac{f(x_n) - f(x_{n-1})}{x_n - x_{n-1}}
\end{equation}

Jika kita substitusikan pendekatan tersebut ke dalam rumus Newton-Raphson, kita akan mendapatkan rumus iterasi Secant:
\begin{equation}
x_{n+1} = x_n - f(x_n) \frac{x_n - x_{n-1}}{f(x_n) - f(x_{n-1})}
\end{equation}

Di mana:
\begin{itemize}
  \item $x_{n+1}$ adalah hampiran akar baru.
  \item $x_n$ dan $x_{n-1}$ adalah dua hampiran akar sebelumnya.
\end{itemize}

\section{Algoritma Iterasi}
\begin{enumerate}
  \item Tentukan dua tebakan awal, $x_0$ dan $x_1$.
  \item Hitung $x_2$ menggunakan rumus iterasi Secant.
  \item Gunakan nilai $x_1$ dan $x_2$ untuk mencari $x_3$, dan seterusnya.
  \item Iterasi berhenti jika $|x_{n+1} - x_n| < \epsilon$ atau $|f(x_{n+1})| < \epsilon$.
\end{enumerate}

\section{Kelebihan dan Kekurangan}
\subsection{Kelebihan}
\begin{itemize}
  \item Tidak memerlukan perhitungan turunan fungsi $f'(x)$, sehingga sangat berguna untuk fungsi yang kompleks.
  \item Kecepatan konvergensi lebih cepat daripada metode Biseksi dan Regula Falsi (mendekati superlinear).
\end{itemize}

\subsection{Kekurangan}
\begin{itemize}
  \item Membutuhkan dua tebakan awal ($x_0$ dan $x_1$).
  \item Karena merupakan metode terbuka, ada kemungkinan iterasi menjadi divergen (menjauh dari akar) jika tebakan awal tidak tepat.
\end{itemize}

\end{document}